\section{证据}
\begin{enumerate}
  \item \zd{能读出命题老师的险恶用心}。太恶心了。\zd{做题先揪出本案事实是什么}。
  \item 理论分类。
        \begin{enumerate}
          \item 证据的来源。原始证据、传来证据（复制、转述）。\al{借条的复印件}。
          \item 证据的内容。跟证据来源、载体、证明力大小\footnote{描述为无法核对、存在疑点、有利害关系。}等等没有关系。\zd{内容能完整证明事实}。
                \begin{enumerate}
                  \item \zd{直接证据}。能够单独、直接、完整证明（\zd{肯定 \hz 否定}）\zd{待证事实}的（\zd{全部，而非部分、片段}）。
                  \item \zd{间接证据}。靠不同的证据证明不同的事实拼接得出结论。\al{拼图；转账和借钱；要约}。
                \end{enumerate}
          \item 证据与证明责任的关系。同一个事实，不同证据，才有本证、反证之分。\al{甲借条证明借钱了，乙收条证明还钱了}。
                \begin{enumerate}
                  \item 本证。\zd{待证事实是什么 \yjt 谁承担证明责任 \yjt 承担的人拿出来 \yjt 本证}。\zd{需达到证明标准}。\al{借条}。
                  \item 反证。无需达到证明标准，只是为了反驳对方主张的事实 \yjt 事实不清证据不明 \yjt 对方倒霉。\al{收条}。
                \end{enumerate}
        \end{enumerate}
  \item 证据的法定分类。8 类。\zd{判断分类，只看原始形态，不看传来形态}。\al{甲的证人提出的书面证词属于书证}。
        \begin{enumerate}
          \item 当事人陈述、勘验笔录没什么考点。
          \item \zd{物证、书证、视听资料、电子数据}。
                \begin{enumerate}
                  \item \zd{内容} \yjt 书证。
                  \item 物品的\zd{物理属性} \yjt 物证。颜色、大小、痕迹、损坏状态等。
                  \item 模拟信号 \yjt 视听资料。\al{录音带}。实践中很少很少很少很少用了，但考试还是考啊。
                  \item 电子信号\yjt 电子数据。储存在\zd{电子介质}当中。\al{短信；通话录音}。
                \end{enumerate}
          \item 书证的规则。同样适用于视听资料电子数据（因为都是用内容证明事实）。
                \begin{enumerate}
                  \item \zd{最佳证据规则}。原则提供原件，有困难复印件。
                  \item 文书真实性规则。
                        \begin{enumerate}
                          \item 公文书证。直接推定为真实，有相反证据足以推翻的除外。
                          \item 私文书证。\zd{由主张以该私文书证证明案件事实的承担证明责任} \yjt 这张“纸“ + 签章是真的。
                                \begin{itemize}
                                  \item \zd{有签名、盖章、捺印，推定真实}（内容真实）\yjt 只需证明签名、盖章等是否伪造（形式真实）。
                                  \item \al{自行委托鉴定}视为私文书证。
                                \end{itemize}
                        \end{enumerate}
                  \item 文书提出命令。书证在对方当事人的控制之下。
                        \begin{enumerate}
                          \item 拒不提交 \yjt 推定书证\zd{内容}为真实。
                          \item 毁损（\al{当庭吞借条}）\yjt 推定主张的案件\zd{事实}为真实。
                        \end{enumerate}
                  \item 域外书证的认证规则。\yi 一般 \yjt 公证 \er 身份关系 \yjt 公证 + 使领馆认证 \san 外文 \yjt 中文译本。
                \end{enumerate}
          \item 证人证言。注意\zd{程序问题}。
                \begin{enumerate}
                  \item 资格。知晓案件事实，能够正确表达。
                  \item 与案件\zd{有利害关系}的可以作为证人，其证言有证明力，只是证明力小\zd{不能单独定案}。
                  \item 出庭。举证期限内申请 \hz 依职权。
                  \item 证人身份冲突，优先适用证人身份。\al{甲既是诉讼代理人，又是目击证人}。
                  \item 不适用回避，必须出庭作证，\zd{否则不能作为定案根据}。
                        \begin{enumerate}
                          \item 视为出庭。庭前的证据交换、法院组织的调查  询问已经出席。
                          \item 法定原因。健康、路途遥远交通不便、自然灾害、不可抗力。
                          \item 议定原因。双方均同意，法院准许。
                        \end{enumerate}
                  \item 出庭。
                        \begin{enumerate}
                          \item 应当签署保证书（无限制民事行为能力人除外）。拒绝 \yjt 不能作证，自行承担费用。
                          \item 如实陈述。不允许使用猜测推断 \hz 评论性的语言（\zd{意见证据规则}）。看到什么，听到什么。
                          \item \zd{证前不得旁听}庭审，不允许以\zd{宣读}（经过排练，不能保证真实性）方式陈述证言。
                          \item 费用败诉方承担，申请方当垫付。
                        \end{enumerate}
                \end{enumerate}
          \item 鉴定意见。
                \begin{enumerate}
                  \item 当事人申请 \hz 法院依职权。
                  \item 鉴定人的确定。\yi 申请的 \yjt 双方协商 \yjt 不成法院指定 \er 法院依职权 \yjt 询问双方意见，直接指定。
                  \item 有权了解进行鉴定所需要的案件材料，可以询问当事人、证人。
                  \item 鉴定前的准备。
                        \begin{enumerate}
                          \item 承诺书。
                          \item 对鉴定材料/检材的质证\footnote{出示证据，双方就证据真实性、关联性、合法性以及证明力的有无和大小问题来发表意见。}，
                                \zd{未经质证的材料不得作为鉴定的根据}。
                        \end{enumerate}
                  \item \zd{不同意见的如实注明，不允许少数服从多少}。
                  \item 当事人有异议的，提出书面异议，通知鉴定人书面答复。出庭费用适用证人。
                  \item 应当出庭拒不出庭的。\al{有正当理由不能出庭的除外 ×}。因为鉴定意见可替代。
                        \begin{enumerate}
                          \item \zd{鉴定意见不允许作为定案依据}。
                          \item 要求返还鉴定费用。
                          \item 法院可以向有关主管部门 \hz 行业组织提出处罚建议。
                        \end{enumerate}
                \end{enumerate}
          \item \zd{有专门知识的人/专家辅助人}。陈述的是研究学习的内容，可替代。\zd{我国不承认专家证人}。
                \begin{enumerate}
                  \item 举证期限内申请（因为是站在当事人立场，所以法院不会依职权去做），1 - 2 名。
                  \item 作用。\yi 对鉴定意见提出意见 \yjt 帮当事人质证 \er 对专业问题提出意见 \yjt 帮当事人陈述。
                  \item 拟制为\zd{当事人陈述}。
                  \item 费用申请方出，不适用回避。
                \end{enumerate}
        \end{enumerate}
\end{enumerate}
